% --- Archon physics formalization source begin ---
% archon:physics
% archon:covers PhyXMiniProblems/problem_phyx_mini_0798.lean
% archon:source-report reports/phyx_mini/problem_phyx_mini_0798.source.json

\chapter{Physics problem phyx\_mini\_0798}
\label{ch:PhyXMiniProblems_problem_phyx_mini_0798}

\paragraph{Problem source.}
\#\# Physical scenario

A waitress shoves a ketchup bottle with mass \textbackslash{}( 0.45 \textbackslash{}, \textbackslash{}text\{kg\} \textbackslash{}) to her right along a smooth, level lunch counter as shown in figure. The bottle leaves her hand moving at \textbackslash{}( 2.0 \textbackslash{}, \textbackslash{}text\{m/s\} \textbackslash{}), then slows down as it slides because of a constant horizontal friction force exerted on it by the countertop. It slides for \textbackslash{}( 1.0 \textbackslash{}, \textbackslash{}text\{m\} \textbackslash{}) before coming to rest.

\#\# Question

What are the magnitude of the friction force acting on the bottle?

\#\# Answer choices

- A: A: 0.45N

- B: B: 0.90N

- C: C: 1.35N

- D: D: 1.80N

\#\# Figure caption (auxiliary; use the image as primary evidence)

The image shows a physics diagram depicting two bottles on a horizontal line. The objects and elements include:

1. **Bottles**: Two bottles are placed on the horizontal line, one at point \textbackslash{}(O\textbackslash{}) and the other at point \textbackslash{}(X\textbackslash{}).

2. **Labels and Text**:
   - The mass of each bottle is labeled as \textbackslash{}(m = 0.45 \textbackslash{}, \textbackslash{}text\{kg\}\textbackslash{}).
   - The initial velocity of the bottle at \textbackslash{}(O\textbackslash{}) is labeled as \textbackslash{}(v\_\{0x\} = 2.0 \textbackslash{}, \textbackslash{}text\{m/s\}\textbackslash{}), with an arrow pointing to the right, indicating the direction of motion.
   - The velocity of the bottle at \textbackslash{}(X\textbackslash{}) is labeled as \textbackslash{}(v\_\{x\} = 0\textbackslash{}).

3. **Distance**: The distance between the two points \textbackslash{}(O\textbackslash{}) and \textbackslash{}(X\textbackslash{}) is marked as \textbackslash{}(1.0 \textbackslash{}, \textbackslash{}text\{m\}\textbackslash{}).

4. **Line**: A horizontal line connects the two points, with \textbackslash{}(O\textbackslash{}) on the left and \textbackslash{}(X\textbackslash{}) on the right.

The setup appears to depict a scenario involving motion along a straight path, likely meant to illustrate principles of kinematics or dynamics.

\#\# Dataset metadata

Dynamics; Physical Model Grounding Reasoning, Multi-Formula Reasoning

\paragraph{Recorded answer/context.}
B: B: 0.90N

\paragraph{Figure/image path.}
/root/proposal\_for\_physic/science-mango/phyx\_mini\_run/phyx\_data/test\_image/798.png

\paragraph{Formalization target.}
create a compiling Lean file with sorry bodies at `PhyXMiniProblems/problem\_phyx\_mini\_0798.lean`. The Lean declarations must preserve the physical quantities, dimensions or dimensional roles, figure labels, governing-law hypotheses, and final relation expressed by this problem.
Use Mathlib/Physlib names found through LeanExplore where available. If a physics API is missing, introduce faithful local abstractions rather than scalar placeholder aliases.

\begin{theorem}[Physics formalization target]
\label{thm:physics:phyx_mini_0798:target}
\lean{PhyXMiniProblems.ProblemPhyXMini0798.frictionForceMagnitude_is_0_90_newtons}
\uses{def:physics:phyx-mini-0798:phyxminiproblems-problemphyxmini0798-forcemagnitudeinnewtons, def:physics:phyx-mini-0798:phyxminiproblems-problemphyxmini0798-bottleslidesetup, def:physics:phyx-mini-0798:phyxminiproblems-problemphyxmini0798-matchesbottleslidescenario, def:physics:phyx-mini-0798:phyxminiproblems-problemphyxmini0798-matchesprimaryfigure, def:physics:phyx-mini-0798:phyxminiproblems-problemphyxmini0798-hasphysicalbottleslideparameters, def:physics:phyx-mini-0798:phyxminiproblems-problemphyxmini0798-satisfiesconstantfrictionworkenergylaws, def:physics:phyx-mini-0798:phyxminiproblems-problemphyxmini0798-recordedanswerchoice, def:physics:phyx-mini-0798:phyxminiproblems-problemphyxmini0798-matchesanswerchoice}
The assigned autoformalize agent should translate this physics problem into Lean declarations in the covered file, with theorem and lemma proof bodies written as `by sorry`.
\end{theorem}
\begin{proof}
This is an autoformalization task, not a proof task. Produce statements that can later be proved by the physics prover without weakening the physical model.
\end{proof}

% --- Archon declaration topology begin ---
\section{Declaration topology}
\begin{definition}[force Dimension]
  \label{def:physics:phyx-mini-0798:phyxminiproblems-problemphyxmini0798-forcedimension}
  \lean{PhyXMiniProblems.ProblemPhyXMini0798.forceDimension}
  Physical force dimension, M L T⁻²
\end{definition}
\begin{proof}
  This is the declared model definition of force Dimension.
\end{proof}
\begin{definition}[energy Dimension]
  \label{def:physics:phyx-mini-0798:phyxminiproblems-problemphyxmini0798-energydimension}
  \lean{PhyXMiniProblems.ProblemPhyXMini0798.energyDimension}
  Physical energy/work dimension, M L² T⁻²
\end{definition}
\begin{proof}
  This is the declared model definition of energy Dimension.
\end{proof}
\begin{definition}[Mass Quantity]
  \label{def:physics:phyx-mini-0798:phyxminiproblems-problemphyxmini0798-massquantity}
  \lean{PhyXMiniProblems.ProblemPhyXMini0798.MassQuantity}
  A nonnegative physical mass
\end{definition}
\begin{proof}
  This is the declared model definition of Mass Quantity.
\end{proof}
\begin{definition}[Horizontal Position Quantity]
  \label{def:physics:phyx-mini-0798:phyxminiproblems-problemphyxmini0798-horizontalpositionquantity}
  \lean{PhyXMiniProblems.ProblemPhyXMini0798.HorizontalPositionQuantity}
  A signed horizontal position
\end{definition}
\begin{proof}
  This is the declared model definition of Horizontal Position Quantity.
\end{proof}
\begin{definition}[Length Magnitude Quantity]
  \label{def:physics:phyx-mini-0798:phyxminiproblems-problemphyxmini0798-lengthmagnitudequantity}
  \lean{PhyXMiniProblems.ProblemPhyXMini0798.LengthMagnitudeQuantity}
  A nonnegative travel distance
\end{definition}
\begin{proof}
  This is the declared model definition of Length Magnitude Quantity.
\end{proof}
\begin{definition}[Horizontal Velocity Quantity]
  \label{def:physics:phyx-mini-0798:phyxminiproblems-problemphyxmini0798-horizontalvelocityquantity}
  \lean{PhyXMiniProblems.ProblemPhyXMini0798.HorizontalVelocityQuantity}
  A signed horizontal velocity, positive to the right
\end{definition}
\begin{proof}
  This is the declared model definition of Horizontal Velocity Quantity.
\end{proof}
\begin{definition}[Force Magnitude Quantity]
  \label{def:physics:phyx-mini-0798:phyxminiproblems-problemphyxmini0798-forcemagnitudequantity}
  \lean{PhyXMiniProblems.ProblemPhyXMini0798.ForceMagnitudeQuantity}
  \uses{def:physics:phyx-mini-0798:phyxminiproblems-problemphyxmini0798-forcedimension}
  A nonnegative force magnitude
\end{definition}
\begin{proof}
  This is the declared model definition of Force Magnitude Quantity.
\end{proof}
\begin{definition}[Horizontal Force Quantity]
  \label{def:physics:phyx-mini-0798:phyxminiproblems-problemphyxmini0798-horizontalforcequantity}
  \lean{PhyXMiniProblems.ProblemPhyXMini0798.HorizontalForceQuantity}
  \uses{def:physics:phyx-mini-0798:phyxminiproblems-problemphyxmini0798-forcedimension}
  A signed horizontal force, positive to the right
\end{definition}
\begin{proof}
  This is the declared model definition of Horizontal Force Quantity.
\end{proof}
\begin{definition}[Energy Magnitude Quantity]
  \label{def:physics:phyx-mini-0798:phyxminiproblems-problemphyxmini0798-energymagnitudequantity}
  \lean{PhyXMiniProblems.ProblemPhyXMini0798.EnergyMagnitudeQuantity}
  \uses{def:physics:phyx-mini-0798:phyxminiproblems-problemphyxmini0798-energydimension}
  A nonnegative kinetic-energy quantity
\end{definition}
\begin{proof}
  This is the declared model definition of Energy Magnitude Quantity.
\end{proof}
\begin{definition}[Signed Work Quantity]
  \label{def:physics:phyx-mini-0798:phyxminiproblems-problemphyxmini0798-signedworkquantity}
  \lean{PhyXMiniProblems.ProblemPhyXMini0798.SignedWorkQuantity}
  \uses{def:physics:phyx-mini-0798:phyxminiproblems-problemphyxmini0798-energydimension}
  Signed work; opposing friction does negative work
\end{definition}
\begin{proof}
  This is the declared model definition of Signed Work Quantity.
\end{proof}
\begin{definition}[nonnegative SIReadout]
  \label{def:physics:phyx-mini-0798:phyxminiproblems-problemphyxmini0798-nonnegativesireadout}
  \lean{PhyXMiniProblems.ProblemPhyXMini0798.nonnegativeSIReadout}
  Read a nonnegative dimensionful quantity in coherent SI units
\end{definition}
\begin{proof}
  This is the declared model definition of nonnegative SIReadout.
\end{proof}
\begin{definition}[signed SIReadout]
  \label{def:physics:phyx-mini-0798:phyxminiproblems-problemphyxmini0798-signedsireadout}
  \lean{PhyXMiniProblems.ProblemPhyXMini0798.signedSIReadout}
  Read a signed dimensionful quantity in coherent SI units
\end{definition}
\begin{proof}
  This is the declared model definition of signed SIReadout.
\end{proof}
\begin{definition}[mass In Kilograms]
  \label{def:physics:phyx-mini-0798:phyxminiproblems-problemphyxmini0798-massinkilograms}
  \lean{PhyXMiniProblems.ProblemPhyXMini0798.massInKilograms}
  \uses{def:physics:phyx-mini-0798:phyxminiproblems-problemphyxmini0798-massquantity, def:physics:phyx-mini-0798:phyxminiproblems-problemphyxmini0798-nonnegativesireadout}
  Kilogram readout of a mass
\end{definition}
\begin{proof}
  This is the declared model definition of mass In Kilograms.
\end{proof}
\begin{definition}[position In Meters]
  \label{def:physics:phyx-mini-0798:phyxminiproblems-problemphyxmini0798-positioninmeters}
  \lean{PhyXMiniProblems.ProblemPhyXMini0798.positionInMeters}
  \uses{def:physics:phyx-mini-0798:phyxminiproblems-problemphyxmini0798-horizontalpositionquantity, def:physics:phyx-mini-0798:phyxminiproblems-problemphyxmini0798-signedsireadout}
  Metre readout of a signed horizontal position
\end{definition}
\begin{proof}
  This is the declared model definition of position In Meters.
\end{proof}
\begin{definition}[length In Meters]
  \label{def:physics:phyx-mini-0798:phyxminiproblems-problemphyxmini0798-lengthinmeters}
  \lean{PhyXMiniProblems.ProblemPhyXMini0798.lengthInMeters}
  \uses{def:physics:phyx-mini-0798:phyxminiproblems-problemphyxmini0798-lengthmagnitudequantity, def:physics:phyx-mini-0798:phyxminiproblems-problemphyxmini0798-nonnegativesireadout}
  Metre readout of a nonnegative length
\end{definition}
\begin{proof}
  This is the declared model definition of length In Meters.
\end{proof}
\begin{definition}[velocity In Meters Per Second]
  \label{def:physics:phyx-mini-0798:phyxminiproblems-problemphyxmini0798-velocityinmeterspersecond}
  \lean{PhyXMiniProblems.ProblemPhyXMini0798.velocityInMetersPerSecond}
  \uses{def:physics:phyx-mini-0798:phyxminiproblems-problemphyxmini0798-horizontalvelocityquantity, def:physics:phyx-mini-0798:phyxminiproblems-problemphyxmini0798-signedsireadout}
  Signed metre-per-second readout of a horizontal velocity
\end{definition}
\begin{proof}
  This is the declared model definition of velocity In Meters Per Second.
\end{proof}
\begin{definition}[force Magnitude In Newtons]
  \label{def:physics:phyx-mini-0798:phyxminiproblems-problemphyxmini0798-forcemagnitudeinnewtons}
  \lean{PhyXMiniProblems.ProblemPhyXMini0798.forceMagnitudeInNewtons}
  \uses{def:physics:phyx-mini-0798:phyxminiproblems-problemphyxmini0798-forcemagnitudequantity, def:physics:phyx-mini-0798:phyxminiproblems-problemphyxmini0798-nonnegativesireadout}
  Newton readout of a force magnitude
\end{definition}
\begin{proof}
  This is the declared model definition of force Magnitude In Newtons.
\end{proof}
\begin{definition}[horizontal Force In Newtons]
  \label{def:physics:phyx-mini-0798:phyxminiproblems-problemphyxmini0798-horizontalforceinnewtons}
  \lean{PhyXMiniProblems.ProblemPhyXMini0798.horizontalForceInNewtons}
  \uses{def:physics:phyx-mini-0798:phyxminiproblems-problemphyxmini0798-horizontalforcequantity, def:physics:phyx-mini-0798:phyxminiproblems-problemphyxmini0798-signedsireadout}
  Signed newton readout of a horizontal force
\end{definition}
\begin{proof}
  This is the declared model definition of horizontal Force In Newtons.
\end{proof}
\begin{definition}[energy In Joules]
  \label{def:physics:phyx-mini-0798:phyxminiproblems-problemphyxmini0798-energyinjoules}
  \lean{PhyXMiniProblems.ProblemPhyXMini0798.energyInJoules}
  \uses{def:physics:phyx-mini-0798:phyxminiproblems-problemphyxmini0798-energymagnitudequantity, def:physics:phyx-mini-0798:phyxminiproblems-problemphyxmini0798-nonnegativesireadout}
  Joule readout of a nonnegative energy
\end{definition}
\begin{proof}
  This is the declared model definition of energy In Joules.
\end{proof}
\begin{definition}[work In Joules]
  \label{def:physics:phyx-mini-0798:phyxminiproblems-problemphyxmini0798-workinjoules}
  \lean{PhyXMiniProblems.ProblemPhyXMini0798.workInJoules}
  \uses{def:physics:phyx-mini-0798:phyxminiproblems-problemphyxmini0798-signedworkquantity, def:physics:phyx-mini-0798:phyxminiproblems-problemphyxmini0798-signedsireadout}
  Signed joule readout of work
\end{definition}
\begin{proof}
  This is the declared model definition of work In Joules.
\end{proof}
\begin{definition}[Figure Point]
  \label{def:physics:phyx-mini-0798:phyxminiproblems-problemphyxmini0798-figurepoint}
  \lean{PhyXMiniProblems.ProblemPhyXMini0798.FigurePoint}
  The two bottle positions explicitly labelled in the supplied image
\end{definition}
\begin{proof}
  This is the declared model definition of Figure Point.
\end{proof}
\begin{definition}[Horizontal Direction]
  \label{def:physics:phyx-mini-0798:phyxminiproblems-problemphyxmini0798-horizontaldirection}
  \lean{PhyXMiniProblems.ProblemPhyXMini0798.HorizontalDirection}
  Horizontal directions distinguished by the problem
\end{definition}
\begin{proof}
  This is the declared model definition of Horizontal Direction.
\end{proof}
\begin{definition}[Counter Geometry]
  \label{def:physics:phyx-mini-0798:phyxminiproblems-problemphyxmini0798-countergeometry}
  \lean{PhyXMiniProblems.ProblemPhyXMini0798.CounterGeometry}
  Geometry of the supporting counter
\end{definition}
\begin{proof}
  This is the declared model definition of Counter Geometry.
\end{proof}
\begin{definition}[Contact Force Model]
  \label{def:physics:phyx-mini-0798:phyxminiproblems-problemphyxmini0798-contactforcemodel}
  \lean{PhyXMiniProblems.ProblemPhyXMini0798.ContactForceModel}
  The idealized contact force during the slide
\end{definition}
\begin{proof}
  This is the declared model definition of Contact Force Model.
\end{proof}
\begin{definition}[Bottle Slide Figure]
  \label{def:physics:phyx-mini-0798:phyxminiproblems-problemphyxmini0798-bottleslidefigure}
  \lean{PhyXMiniProblems.ProblemPhyXMini0798.BottleSlideFigure}
  \uses{def:physics:phyx-mini-0798:phyxminiproblems-problemphyxmini0798-figurepoint, def:physics:phyx-mini-0798:phyxminiproblems-problemphyxmini0798-horizontaldirection}
  Literal diagram information transcribed from the primary image
\end{definition}
\begin{proof}
  This is the declared model definition of Bottle Slide Figure.
\end{proof}
\begin{definition}[Bottle Slide Setup]
  \label{def:physics:phyx-mini-0798:phyxminiproblems-problemphyxmini0798-bottleslidesetup}
  \lean{PhyXMiniProblems.ProblemPhyXMini0798.BottleSlideSetup}
  \uses{def:physics:phyx-mini-0798:phyxminiproblems-problemphyxmini0798-massquantity, def:physics:phyx-mini-0798:phyxminiproblems-problemphyxmini0798-horizontalpositionquantity, def:physics:phyx-mini-0798:phyxminiproblems-problemphyxmini0798-lengthmagnitudequantity, def:physics:phyx-mini-0798:phyxminiproblems-problemphyxmini0798-horizontalvelocityquantity, def:physics:phyx-mini-0798:phyxminiproblems-problemphyxmini0798-forcemagnitudequantity, def:physics:phyx-mini-0798:phyxminiproblems-problemphyxmini0798-horizontalforcequantity, def:physics:phyx-mini-0798:phyxminiproblems-problemphyxmini0798-energymagnitudequantity, def:physics:phyx-mini-0798:phyxminiproblems-problemphyxmini0798-signedworkquantity, def:physics:phyx-mini-0798:phyxminiproblems-problemphyxmini0798-figurepoint, def:physics:phyx-mini-0798:phyxminiproblems-problemphyxmini0798-horizontaldirection, def:physics:phyx-mini-0798:phyxminiproblems-problemphyxmini0798-countergeometry, def:physics:phyx-mini-0798:phyxminiproblems-problemphyxmini0798-contactforcemodel, def:physics:phyx-mini-0798:phyxminiproblems-problemphyxmini0798-bottleslidefigure}
  The physical bottle and its independent dynamical observables. The force, kinetic-energy, and work fields are connected only by the governing laws below; no field is assigned the requested numerical answer
\end{definition}
\begin{proof}
  This is the declared model definition of Bottle Slide Setup.
\end{proof}
\begin{definition}[Matches Bottle Slide Scenario]
  \label{def:physics:phyx-mini-0798:phyxminiproblems-problemphyxmini0798-matchesbottleslidescenario}
  \lean{PhyXMiniProblems.ProblemPhyXMini0798.MatchesBottleSlideScenario}
  \uses{def:physics:phyx-mini-0798:phyxminiproblems-problemphyxmini0798-bottleslidesetup}
  Qualitative physical idealizations stated in the problem prose
\end{definition}
\begin{proof}
  This is the declared model definition of Matches Bottle Slide Scenario.
\end{proof}
\begin{definition}[Matches Primary Figure]
  \label{def:physics:phyx-mini-0798:phyxminiproblems-problemphyxmini0798-matchesprimaryfigure}
  \lean{PhyXMiniProblems.ProblemPhyXMini0798.MatchesPrimaryFigure}
  \uses{def:physics:phyx-mini-0798:phyxminiproblems-problemphyxmini0798-massinkilograms, def:physics:phyx-mini-0798:phyxminiproblems-problemphyxmini0798-lengthinmeters, def:physics:phyx-mini-0798:phyxminiproblems-problemphyxmini0798-velocityinmeterspersecond, def:physics:phyx-mini-0798:phyxminiproblems-problemphyxmini0798-figurepoint, def:physics:phyx-mini-0798:phyxminiproblems-problemphyxmini0798-bottleslidesetup}
  Objects, labels, geometry, and numerical inscriptions visible in image 798. The physical-quantity equalities connect those inscriptions to the setup but contain no assertion about the friction magnitude
\end{definition}
\begin{proof}
  This is the declared model definition of Matches Primary Figure.
\end{proof}
\begin{definition}[Has Physical Bottle Slide Parameters]
  \label{def:physics:phyx-mini-0798:phyxminiproblems-problemphyxmini0798-hasphysicalbottleslideparameters}
  \lean{PhyXMiniProblems.ProblemPhyXMini0798.HasPhysicalBottleSlideParameters}
  \uses{def:physics:phyx-mini-0798:phyxminiproblems-problemphyxmini0798-massinkilograms, def:physics:phyx-mini-0798:phyxminiproblems-problemphyxmini0798-positioninmeters, def:physics:phyx-mini-0798:phyxminiproblems-problemphyxmini0798-lengthinmeters, def:physics:phyx-mini-0798:phyxminiproblems-problemphyxmini0798-velocityinmeterspersecond, def:physics:phyx-mini-0798:phyxminiproblems-problemphyxmini0798-bottleslidesetup}
  Positivity and endpoint ordering selecting the physical branch
\end{definition}
\begin{proof}
  This is the declared model definition of Has Physical Bottle Slide Parameters.
\end{proof}
\begin{definition}[Satisfies Constant Friction Work Energy Laws]
  \label{def:physics:phyx-mini-0798:phyxminiproblems-problemphyxmini0798-satisfiesconstantfrictionworkenergylaws}
  \lean{PhyXMiniProblems.ProblemPhyXMini0798.SatisfiesConstantFrictionWorkEnergyLaws}
  \uses{def:physics:phyx-mini-0798:phyxminiproblems-problemphyxmini0798-horizontalpositionquantity, def:physics:phyx-mini-0798:phyxminiproblems-problemphyxmini0798-massinkilograms, def:physics:phyx-mini-0798:phyxminiproblems-problemphyxmini0798-positioninmeters, def:physics:phyx-mini-0798:phyxminiproblems-problemphyxmini0798-lengthinmeters, def:physics:phyx-mini-0798:phyxminiproblems-problemphyxmini0798-velocityinmeterspersecond, def:physics:phyx-mini-0798:phyxminiproblems-problemphyxmini0798-forcemagnitudeinnewtons, def:physics:phyx-mini-0798:phyxminiproblems-problemphyxmini0798-horizontalforceinnewtons, def:physics:phyx-mini-0798:phyxminiproblems-problemphyxmini0798-energyinjoules, def:physics:phyx-mini-0798:phyxminiproblems-problemphyxmini0798-workinjoules, def:physics:phyx-mini-0798:phyxminiproblems-problemphyxmini0798-figurepoint, def:physics:phyx-mini-0798:phyxminiproblems-problemphyxmini0798-bottleslidesetup}
  The scalar laws use rightward as positive. The signed friction field is therefore the negative of its nonnegative magnitude at every position between O and X. The remaining fields state the endpoint distance, translational kinetic energy, work of a constant opposing force, and the work--energy theorem. None states a numerical value for the friction force
\end{definition}
\begin{proof}
  This is the declared model definition of Satisfies Constant Friction Work Energy Laws.
\end{proof}
\begin{lemma}[friction Force Magnitude formula]
  \label{lem:physics:phyx-mini-0798:phyxminiproblems-problemphyxmini0798-frictionforcemagnitude-formula}
  \lean{PhyXMiniProblems.ProblemPhyXMini0798.frictionForceMagnitude_formula}
  \uses{def:physics:phyx-mini-0798:phyxminiproblems-problemphyxmini0798-massinkilograms, def:physics:phyx-mini-0798:phyxminiproblems-problemphyxmini0798-lengthinmeters, def:physics:phyx-mini-0798:phyxminiproblems-problemphyxmini0798-velocityinmeterspersecond, def:physics:phyx-mini-0798:phyxminiproblems-problemphyxmini0798-forcemagnitudeinnewtons, def:physics:phyx-mini-0798:phyxminiproblems-problemphyxmini0798-bottleslidesetup, def:physics:phyx-mini-0798:phyxminiproblems-problemphyxmini0798-hasphysicalbottleslideparameters, def:physics:phyx-mini-0798:phyxminiproblems-problemphyxmini0798-satisfiesconstantfrictionworkenergylaws}
  Eliminating the endpoint kinetic energies and friction work yields the usual constant-friction stopping formula. It is a derived conclusion, not a field or premise of the physical model
\end{lemma}
\begin{proof}
  The conclusion follows from the governing relations and figure readouts named in the statement of friction Force Magnitude formula.
\end{proof}
\begin{definition}[Answer Choice]
  \label{def:physics:phyx-mini-0798:phyxminiproblems-problemphyxmini0798-answerchoice}
  \lean{PhyXMiniProblems.ProblemPhyXMini0798.AnswerChoice}
  Labels of the four displayed multiple-choice answers
\end{definition}
\begin{proof}
  This is the declared model definition of Answer Choice.
\end{proof}
\begin{definition}[force Magnitude Newtons]
  \label{def:physics:phyx-mini-0798:phyxminiproblems-problemphyxmini0798-answerchoice-forcemagnitudenewtons}
  \lean{PhyXMiniProblems.ProblemPhyXMini0798.AnswerChoice.forceMagnitudeNewtons}
  \uses{def:physics:phyx-mini-0798:phyxminiproblems-problemphyxmini0798-answerchoice}
  Force magnitude printed beside each answer choice, in newtons
\end{definition}
\begin{proof}
  This is the declared model definition of force Magnitude Newtons.
\end{proof}
\begin{definition}[recorded Answer Choice]
  \label{def:physics:phyx-mini-0798:phyxminiproblems-problemphyxmini0798-recordedanswerchoice}
  \lean{PhyXMiniProblems.ProblemPhyXMini0798.recordedAnswerChoice}
  \uses{def:physics:phyx-mini-0798:phyxminiproblems-problemphyxmini0798-answerchoice}
  The answer label recorded in the supplied dataset metadata
\end{definition}
\begin{proof}
  This is the declared model definition of recorded Answer Choice.
\end{proof}
\begin{definition}[Matches Answer Choice]
  \label{def:physics:phyx-mini-0798:phyxminiproblems-problemphyxmini0798-matchesanswerchoice}
  \lean{PhyXMiniProblems.ProblemPhyXMini0798.MatchesAnswerChoice}
  \uses{def:physics:phyx-mini-0798:phyxminiproblems-problemphyxmini0798-forcemagnitudequantity, def:physics:phyx-mini-0798:phyxminiproblems-problemphyxmini0798-forcemagnitudeinnewtons, def:physics:phyx-mini-0798:phyxminiproblems-problemphyxmini0798-answerchoice}
  A physical friction magnitude agrees with a displayed answer choice
\end{definition}
\begin{proof}
  This is the declared model definition of Matches Answer Choice.
\end{proof}
% --- Archon declaration topology end ---
% --- Archon physics formalization source end ---
